\vspace{.5cm}
\begin{table}[H]
\caption{Welfare Gains from Each Policy}
\label{tab:welfare}
\centering
\begin{threeparttable}
    \begin{tabular*}{\textwidth}{l @{\extracolsep{\fill}} ccccc}
        \toprule
        \multicolumn{1}{l}{Policy} & \multicolumn{1}{c}{Welfare gain} & \multicolumn{4}{c}{Share of customers in each arm} \\
        \cmidrule{3-6}
        & & $(U,U)$ & $(T,U)$ & $(U,T)$ & $(T,T)$\\
        \midrule
        %\addlinespace
        100\% $(U,U)$ & $0.0$ & $100.0\%$ & $0.0\%$ & $0.0\%$ & $0.0\%$\\
        & (0.0) &&& \\
        \addlinespace
        100\% $(T,U)$ & $311.8$ & $0.0\%$ & $100.0\%$ & $0.0\%$ & $0.0\%$\\
        & $(378.4)$ &&& \\
        \addlinespace
        100\% $(U,T)$ & $470.8$ & $0.0\%$ & $0.0\%$ & $100.0\%$ & $0.0\%$\\
        & $(457.5)$ &&& \\
        \addlinespace
        100\% $(T,T)$ & $463.9$ & $0.0\%$ & $0.0\%$ & $0.0\%$ & $100.0\%$\\
        & $(452.2)$ &&& \\
        \addlinespace
        Static targeting I ($\pi_{S(I)}^{\ast}$) & $770.6$ & $45.6\%$ & $0.0\%$ & $0.0\%$ & $54.4\%$\\
        & $(283.7)$ &&& \\
        \addlinespace
        Static targeting II ($\pi_{S(\two)}^{\ast}$) & $845.3$ & $3.1\%$ & $31.3\%$ & $41.5\%$ & $24.0\%$\\
        & $(348.9)$ &&& \\
        \addlinespace
        Dynamic targeting ($\pi^\ast$) & $1684.3$ & $19.5\%$ & $22.9\%$ & $25.6\%$ & $32.0\%$\\
        & $(303.1)$ &&& \\
        %\addlinespace
        \bottomrule
    \end{tabular*}
\begin{tablenotes}[flushleft]
\footnotesize\item Notes: This table summarizes characteristics of four benchmark policies (100\% $(U,U)$, 100\% $(T,U)$, 100\% $(U,T)$, and 100\% $(T,T)$), static targeting (I) ($\pi_{S(I)}^{\ast}$), static targeting (II) ($\pi_{S(\two)}^{\ast}$), and dynamic targeting ($\pi^{\ast}$). The column titled ``Welfare Gain'' shows the estimated welfare gains in JPY per household, with these standard errors in parentheses. The monetary unit is given as 1 ¢ = 1 JPY, the exchange rate in the summer of 2020. %All of these estimates are obtained from the test data. The columns grouped by label ``Share of customers in each term'' give the estimated share of each arm.
\end{tablenotes}
\end{threeparttable}
\end{table}
\vspace{.5cm}